\documentclass[11pt,a4paper]{article}
\usepackage[T1]{fontenc}
\usepackage{lmodern,amsmath,amssymb,amsthm,microtype,booktabs}
\usepackage[margin=25mm]{geometry}
\usepackage[colorlinks=true,urlcolor=blue,linkcolor=blue]{hyperref}
\newtheorem{theorem}{Theorem}
\newtheorem{lemma}{Lemma}
\newcommand{\dd}{\mathop{}\!\mathrm d}
\newcommand{\Q}{Q_\nu}
\title{A full fixed-\(K\) Barnes--Bessel limit\\with an integrated first correction}
\author{Alan Day \and Chengheng\thanks{AI research assistant (OpenAI).}}
\date{Working research note --- 30 October 2026}
\begin{document}\maketitle
\begin{abstract}
An exact unit-shift recurrence controls the positive horizontal tail
of the recoil Barnes kernel uniformly in the mass. A complementary
vertical-tail estimate permits a full fixed-\(K,\alpha\) limit to the
two-branch Bessel integral. A logarithmically growing truncation then
justifies the first integrated correction with an absolute
\(O(\nu^{-2})\) remainder. Constants depend on \(K,\alpha\).
This addresses the global gap in the preceding note, without asserting
a joint growing-\(K\) limit or a physical finite-\(G\) four-point function.
\end{abstract}
\section{Object and fixed-parameter statement}
Set the de Sitter radius to one, \(u=v=0\), \(K>0\), real \(\alpha\),
\(\omega=i\alpha\), and \(M=\nu^2/(2K)\).
Our object is the internal chiral Barnes amplitude after the earlier
fixed-mass \(\kappa\to\infty\) limit. Physical state, algebra,
external insertions and normalization remain as yet unidentified.
Define \(t=z+i(M+\nu)\), \(v_0=\nu/2\), and
\[
 b=\tfrac12-i(M-\nu),\quad c=\tfrac12-i(M+\nu),\quad
 C_{M,\nu}=e^{-\pi(M+\nu)}\frac{\Gamma(-iv_0)^4\Gamma(b)}{\Gamma(c)}.
\]
The normalized exact kernel is
\[
 \Q(t)=e^{-i\pi t}\frac{\Gamma(t-iv_0)^4}{\Gamma(-iv_0)^4}
 \frac{\Gamma(b-t)}{\Gamma(b)}\frac{\Gamma(c)}{\Gamma(c+t)}
 \Gamma(\tfrac12-t+\alpha/2)\Gamma(\tfrac12-t-\alpha/2).
 \tag{1}
\]
Fix \(\eta>0\), and an integer \(N>|\alpha|/2+1/2\).
Let \(\mathcal C\) run from \(-N-i\infty\) upward to \(-N+i\eta\),
then right to \(+\infty+i\eta\). For sufficiently large \(\nu\),
this is the pole-free deformation of the original Barnes contour.
Its left ray points upward. Put
\[
 I_\nu=\frac1{2\pi i}\int_{\mathcal C}\Q(t)\,\dd t
       =\frac{\mathfrak B}{C_{M,\nu}},\qquad
 J=-K\bigl[e^{i\pi\alpha}K_\alpha(K)+i\pi I_\alpha(K)\bigr].
\]
\begin{theorem}
For fixed \(K>0\), real \(\alpha\), as \(\nu\to\infty\),
\[
 I_\nu=J\left[1+\frac{i}{\nu}(K^2+\alpha^2-1)\right]
       +O_{K,\alpha}(\nu^{-2}).
 \tag{2}
\]
The error is absolute, including when \(J=0\).
\end{theorem}

\newpage
\section{The exact recurrence gives a uniform horizontal tail}
Gamma recurrence \cite{gamma}, applied directly to (1), yields
\[
 \frac{\Q(t+1)}{\Q(t)}
 =R_\nu(t):=
 \frac{(t-i\nu/2)^4}
 {[(t+\tfrac12-i\nu)^2+M^2][(t+\tfrac12)^2-\alpha^2/4]} .
 \tag{3}
\]
The minus sign from \(e^{-i\pi(t+1)}\) cancels the sign from
\(\Gamma(b-t-1)/\Gamma(b-t)\). This avoids using a fixed-mass
Stirling constant as though it were mass-uniform.

Set
\[
 d=\max(\alpha^2/4-\eta^2,0),\quad A=\max(2,2K,2d),\quad
 a=\nu/2-\eta,\quad L_\nu=M-\nu+\eta .
\]
For \(\nu\ge\max(8K,4\eta,2)\), \(a\ge0\), \(L_\nu\ge M/2\).
For \(t=x+i\eta\), \(q=x+\tfrac12\), \(x\ge A\),
\begin{align*}
 |(t+\tfrac12-i\nu)^2+M^2|&\ge q^2+L_\nu^2,\\
 |(t+\tfrac12)^2-\alpha^2/4|&\ge q^2-d.
\end{align*}
The first inequality factors the expression into
\(q+i(\eta-\nu+M)\) and \(q+i(\eta-\nu-M)\);
the second is the reverse triangle inequality.
Furthermore,
\[
 (1+\tfrac1{2A})^2(2a^2+a^4/A^2)
 \le\frac{75}{128}\frac{\nu^4}{16K^2}
 \le L_\nu^2. \tag{4}
\]
Indeed, after division by \(\nu^4/(16K^2)\), the upper bound is
\((5/4)^2(8K^2/\nu^2+K^2/A^2)\le75/128\).
Since \(q/x\le1+1/(2A)\), (4) implies
\[
 (x^2+a^2)^2q^2\le x^4(q^2+L_\nu^2).
\]
Using \(x\ge2d\), we conclude
\begin{align}
 |R_\nu(x+i\eta)|
 &\le\frac{x^4}{q^2(q^2-d)}
 \le\frac1{(1+1/x)(1+1/(2x))}
 \le\left(\frac{x}{x+1}\right)^{3/2}. \tag{5}
\end{align}
The last inequality follows from \((1+u/2)^2\ge1+u\).

For any real \(R\ge A\), define
\(H_R=\sup_{R\le x\le R+1}|\Q(x+i\eta)|\).
Iteration on each unit lattice, starting at \(x_0\in[R,R+1)\), gives
\[
 |\Q(x+i\eta)|\le H_R\left(\frac{R+1}{x}\right)^{3/2}
 \quad(x\ge R),\qquad
 \int_R^\infty|\Q(x+i\eta)|\,\dd x\le6RH_R. \tag{6}
\]
At fixed \(R=A\), compact convergence below bounds \(H_A\)
uniformly for all sufficiently large \(\nu\). This supplies an
integrable, mass-independent positive-tail majorant.

\newpage
\section{The downward tail and dominated convergence}
Let \(s\ge0\) on \(t=-N-is\). Fixed-real-part two-sided Gamma
bounds follow from Stirling asymptotics and compactness
\cite{stirling}. Write \(B_0=M-\nu\), \(C_0=M+\nu\).
The four large numerator factors, the two large-argument ratios,
and the two remaining Gamma factors bound (1) by a constant times
\[
 \frac{v_0^2}{(v_0+s)^{4N+2}}
 (1+|B_0-s|)^N(C_0+s)^N(1+s)^{2N}
 \exp\!\left[-\frac{7\pi s}{2}
 -\frac{\pi}{2}(|B_0-s|-B_0)\right]. \tag{7}
\]
This remains valid across \(s=B_0\), where the real part of the
potentially small Gamma argument is \(N+1/2>0\).
The exponential is \(e^{-3\pi s}\) for \(s\le B_0\), and
\(e^{-4\pi s+\pi B_0}\le e^{-3\pi s}\) afterwards.

For \(v_0\ge1\), \(M\ge4\nu\), use
\[
 \frac{M+s+1}{(v_0+s)^2}\le\frac{M}{v_0^2}+1=\frac2K+1,
 \qquad \frac{v_0^2}{(v_0+s)^2}\le1.
\]
Also \(1+|B_0-s|\le M+s+1\) and \(C_0+s\le2(M+s+1)\).
Consequently
\[
 |\Q(-N-is)|\le C_{K,\alpha,N}(1+s)^{2N}e^{-3\pi s}.
 \tag{8}
\]
No \(\nu\)-dependent polynomial prefactor remains.
The short segment \(s\in[-\eta,0]\) is compact and harmless.

At fixed \(t\) on \(\mathcal C\), Gamma-ratio asymptotics give
\[
 \Q(t)\longrightarrow
 U(t):=e^{-2\pi it}(K/2)^{2t}
 \Gamma(\tfrac12-t+\alpha/2)\Gamma(\tfrac12-t-\alpha/2).
 \tag{9}
\]
The segment from \(-N+i\eta\) to \(A+i\eta\) is compact;
(6) and (8) dominate the two tails. Thus
\[
 I_\nu\longrightarrow\frac1{2\pi i}\int_{\mathcal C}U(t)\,\dd t=J.
 \tag{10}
\]
The last equality is the exact auxiliary identity derived in the
preceding note using both pole families and the standard Bessel
connection \cite{bessel,continuation}. At integer \(\alpha\),
continue their complete sum, not the separate singular coefficients.
Equation (10) now concerns the full normalized Barnes integral.
It is not merely compact-kernel convergence.

\newpage
\section{A growing cutoff justifies the integrated correction}
Take \(R=\log\nu\), eventually larger than \(A,N,\eta\).
On the truncated contour with \(x\le R+1\), \(s\le R\), the
large Gamma arguments remain in fixed safe sectors. A uniform
Gamma-ratio expansion for shifts \(h=O(\log\nu)\) follows by
Taylor expansion of \(\log\Gamma\), using Stirling derivative
bounds on discs of radius \(O(\log\nu)\):
\[
 \log\frac{\Gamma(z+h)}{\Gamma(z)}
 =h\log z+\frac{h(h-1)}{2z}
   +O\!\left(\frac{|h|+|h|^3}{|z|^2}\right).
 \tag{11}
\]
Log ratios are continued from zero shift. The discs remain a fixed
angular distance from the negative-real cut.
Applying (11) at \(-i\nu/2,b,c\), with shifts \(t,-t,t\),
respectively, gives
\[
 \log\frac{\Q(t)}{U(t)}
 =\frac{4i}{\nu}(t^2-t)
 +O_{K,\alpha}\!\left(\frac{(1+|t|)^3}{\nu^2}\right). \tag{12}
\]
The leading powers contribute only \(O(|t|/\nu^2)\) beyond (9):
\(\log b+\log c=2\log M-i\pi+O(\nu^2/M^2+1/M)\).
The \(b,c\) quadratic terms are \(O(|t|^2/M)\).
Exponentiating (12), uniformly on this growing truncation,
\[
 \left|\Q(t)-U(t)\left[1+\frac{4i}{\nu}(t^2-t)\right]\right|
 \le \frac{C_{K,\alpha}}{\nu^2}(1+|t|)^4|U(t)|. \tag{13}
\]
All polynomial moments of \(|U|\) are integrable on \(\mathcal C\).

It remains to verify that the discarded exact tails are small at
this \emph{moving} cutoff. Equation (8) gives
\[
 \int_R^\infty|\Q(-N-is)|\,\dd s
 \le C(1+R)^{2N}e^{-3\pi R}=o(\nu^{-2}).
\]
On \([R,R+1]+i\eta\), (12) bounds \(|\Q|\le C|U|\).
Reflection and Stirling for \(U\) give
\[
 H_R\le C\exp[-2R\log R+C_KR+O(\log R)] .
\]
Together with (6), the horizontal exact tail is \(o(\nu^{-m})\)
for every fixed \(m>0\). The auxiliary tails with the polynomial
correction enjoy the same required estimates.
Thus integration of (13) yields
\[
 I_\nu=J+\frac{4i}{\nu}\frac1{2\pi i}
 \int_{\mathcal C}(t^2-t)U(t)\,\dd t+O_{K,\alpha}(\nu^{-2}).
\]
For \(\mathcal D=(K/2)\partial_K\) at \emph{fixed \(\alpha\)},
the auxiliary Bessel equation gives
\[
 (\mathcal D^2-\mathcal D)J
 =\frac{K^2+\alpha^2-1}{4}J.
\]
This proves (2). The moving cutoff is essential: the exact kernel
need not have finite polynomially weighted moments on its full
horizontal tail.

\newpage
\section{Normalized ratio and finite numerical checks}
At \(\alpha=0\), \(J(K,0)=-K[K_0(K)+i\pi I_0(K)]\ne0\) for \(K>0\).
Since \(C_{M,\nu}\) is independent of \(\alpha\), the theorem implies
\[
 \frac{\mathfrak B(K,\alpha)}{\mathfrak B(K,0)}
 =\frac{J(K,\alpha)}{J(K,0)}e^{i\alpha^2/\nu}
 +O_{K,\alpha}(\nu^{-2}). \tag{14}
\]
This is an absolute normalized-ratio error. A relative formulation
also needs \(J(K,\alpha)\ne0\).
The exponential agrees with the proved first-order term;
no all-orders resummation is asserted.

We independently integrate (1) on the fixed corner contour.
For \(I_\nu^{(1)}=J[1+i(K^2+\alpha^2-1)/\nu]\), the results are
\begin{center}\begin{tabular}{rrrrr}\toprule
\(K\)&\(\alpha\)&\(\nu\)&\(|I_\nu/I_\nu^{(1)}-1|\)&
\(\nu^2|I_\nu-I_\nu^{(1)}|\)\\\midrule
2&0.7&64&\(3.97586\,10^{-4}\)&19.59115\\
2&0.7&128&\(9.56063\,10^{-5}\)&18.82311\\
2&0.7&256&\(2.37028\,10^{-5}\)&18.66140\\
4&8&128&0.17119939&75063.29\\
4&8&256&0.04863051&75956.09\\
4&8&512&0.01261136&76178.45\\\bottomrule
\end{tabular}\end{center}
The large difference between the last-column scales illustrates why
a fixed-parameter theorem cannot silently become a joint-\(K\) estimate.

The 69 diagnostics comprise 12 exact recurrence checks, 6 mass
inequality checks, 36 rational tail-step samples, 6 quadrature
refinements, 4 precision/combined contour-extent comparisons,
4 empirical decreases of the first-correction error, and
one normalization comparison with the previous original-\(z\) contour.
Six compact-expansion residuals are also recorded as data.
Quadrature uses 18/30 Gauss nodes per segment and 50/65 digits;
the largest mass in each family is checked at 85 digits and
with \((X,T,\eta)=(36,18,1)\) changed to \((44,24,0.75)\).
These are empirical finite-point checks, not interval certification.
The tail proof is the analytic argument (3)--(8).

\section{What remains}
The previous full-integral gap is closed for fixed \(K,\alpha\) on
the symmetric \(u=v=0\), imaginary-\(\omega\) slice.
The constants in this proof depend on \(K,\alpha\);
in particular \(A\) grows like \(\alpha^2\) and \(N\) like \(|\alpha|\).
The next task is explicit uniform control when \(\alpha=Ky\) and
\(K\to\infty\), including branch cancellation and the transverse
variables. The physical external cycles, normalization, disconnected
baseline, finite-\(G\) limit and microscopic holographic dictionary
remain open. This is an author self-reviewed proof draft.

\begin{thebibliography}{9}
\bibitem{gamma} NIST DLMF, \S5.5, Gamma recurrence and reflection,
\url{https://dlmf.nist.gov/5.5}.
\bibitem{stirling} NIST DLMF, \S5.11, asymptotic expansions and remainders,
\url{https://dlmf.nist.gov/5.11}.
\bibitem{bessel} NIST DLMF, \S10.25 and \S10.27,
\url{https://dlmf.nist.gov/10.27}.
\bibitem{continuation} NIST DLMF, 10.34.2, continuation below the cut,
\url{https://dlmf.nist.gov/10.34}.
\end{thebibliography}
\end{document}
